\scanpage{67}
\mspage{16}

. Mais il est bien possible qu'il faille renforcer un peu $W(1)$ \inserted{(saturation)}, en exigeant que \deleted{\unsure{tout}} rétracté d'un $f\in W$ soit dans $W$.

\uline{\unsure{Remarques|Remarque}} Considérons, \inserted{pour $n\in\mathbb N$,} l'\unsure{ens} $W_n\subset\mathrm{Fl}(\mathrm{Cat})$ des ``$n$-équivalences \deleted{\unsure{d'homotopie}} \inserted{d'homotopie}'' de $\mathrm{Cat}$ : \[\begin{aligned}f\in W_n\iff{}&\pi_0(f)\text{ bijectif, et }\pi_i(f,x):\pi_i(X,x)\longrightarrow\pi_i(Y,fx)\\ &\text{bijectif pour }0\unsure{<|\leq}i\leq n.\end{aligned}\] Je croyais que ça \deleted{\unsure{fait}} satisfaisait à \deleted{\unsure{tout}} ce qu'on peut demander à un localiseur, sans avoir pris la peine de l'écrire avec \unsure{soin}. \deleted{\unsure{Bien sûr, ça sat'}} \inserted{Mais ça ne peut} \deleted{\unsure{pour plaire}} \inserted{satisfaire} à $W(4)$ \unsure{jusqu'à|puisqu'à} $W(4\,\mathrm b'')$, puisque on n'a pas $W_n\subset W_\infty$ !

$W(1)$ est trivial sur la définition (et on a sans bavure même la \uline{saturation forte} de $W_i$, i.e. si $f$ \deleted{\unsure{déjà}} \unsure{$\in\mathrm{Cat}$|t.q.} est iso dans $\mathrm{Hot}_{W_i}=(\mathrm{Cat})W_i^{-1}$, alors $f\in W_i$). Je crois me rappeler avoir vérifié $W(2)$, par la même méthode que pour $W_\infty$, en utilisant un critère cohomologique, à la Artin-Mazur, pour \deleted{$f\in$} exprimer $f\in W_i$, et en utilisant les suites spectrales de Leray pour $X\longrightarrow S$ et $Y\longrightarrow S$. La condition $W(3)$ est trivialement satisfaite puisque $W_n\supset W_\infty$ et qu'elle l'est pour $W_\infty$. Donc ça coince là où je m'y attendais le moins : l'axiome des fibrations. Prenons la forme faible : Soit $f:X\longrightarrow Y$ une $W_\omega$-fibration, \margin{$\omega>\infty$\\ $W_\omega\subset W_\infty\subset W_n$} et $f\in W_i$, les fibres\inserted{$X_s$} sont elles $W_i$-asphériques ? SPS $X$ et $Y$ 0-connexes, soit $x\in X_s$, la suite exacte d'homotopie donne \diagram{\begin{tikzcd}[ampersand replacement=\&,column sep=small] \pi_{i+1}(X)\arrow[r, "{\text{iso si }i\leq n-1}"']\&\pi_{i+1}(Y)\arrow[r]\&\pi_i(X_s)\arrow[r]\&\pi_i(X)\arrow[r,"{\mathrm{iso}}","{\text{si }i\leq n}"']\&\pi_i(Y) \end{tikzcd}} donc \[\pi_i(X_s)=1\qquad\text{si }0\leq i\leq n-1,\] mais pour $i=n$, on trouve

% archive sheet number: 67
